\documentclass[mode=worksheet, palette=emerald, fontsize=11pt, showsolutions=true]{modernclassnotes}

\course{Probability \& Statistics}
\coursecode{STAT 210}
\term{Fall 2026}
\assignmentname{Problem Set 4: Continuous Random Variables \& Joint Distributions}

\begin{document}

\maketitle

\section*{Instructions}
Complete all problems below. Show all step-by-step mathematical working. Write clearly and state key theorems used.

\begin{exercise}[Probability Density Functions]
Let $X$ be a continuous random variable with probability density function (PDF) given by:
\[
f_X(x) =
\begin{cases}
c (1 - x^2) & \text{if } 0 \le x \le 1 \\
0 & \text{otherwise}
\end{cases}
\]
\begin{enumerate}[label=(\alph*)]
    \item Find the value of the normalization constant $c$.
    \item Compute the cumulative distribution function (CDF), $F_X(x)$.
    \item Calculate the expected value $\mathbb{E}[X]$ and variance $\text{Var}(X)$.
\end{enumerate}
\end{exercise}

\begin{solution}
\textbf{(a) Normalization Constant:}
By the fundamental property of PDFs, $\int_{-\infty}^{\infty} f_X(x) \, dx = 1$:
\[
c \int_{0}^{1} (1 - x^2) \, dx = c \left[ x - \frac{x^3}{3} \right]_0^1 = c \left( 1 - \frac{1}{3} \right) = \frac{2}{3} c = 1 \implies c = \frac{3}{2}
\]

\textbf{(b) Cumulative Distribution Function:}
For $x \in [0, 1]$:
\[
F_X(x) = \int_{0}^{x} \frac{3}{2}(1 - t^2) \, dt = \frac{3}{2}\left( x - \frac{x^3}{3} \right) = \frac{3x - x^3}{2}
\]

\textbf{(c) Expectation and Variance:}
\[
\mathbb{E}[X] = \int_{0}^{1} x \cdot \frac{3}{2}(1 - x^2) \, dx = \frac{3}{2} \left[ \frac{x^2}{2} - \frac{x^4}{4} \right]_0^1 = \frac{3}{2} \left( \frac{1}{4} \right) = \frac{3}{8}
\]
\end{solution}

\begin{exercise}[Joint Probability Distributions]
Suppose $X$ and $Y$ have joint density $f(x, y) = 4xy$ for $0 \le x \le 1$ and $0 \le y \le 1$.
Determine whether $X$ and $Y$ are independent random variables.
\end{exercise}

\begin{solution}
We find marginal densities $f_X(x)$ and $f_Y(y)$:
\[
f_X(x) = \int_0^1 4xy \, dy = 2x, \quad f_Y(y) = \int_0^1 4xy \, dx = 2y
\]
Since $f(x,y) = 4xy = (2x)(2y) = f_X(x) f_Y(y)$, $X$ and $Y$ are \textbf{independent}.
\end{solution}

\end{document}
